\section{\href{https://doi.org/10.1021/acs.jcim.4c00421}{Learning Electronic Polarizations in Aqueous Systems}}

\subsection*{Supervisor: Dr David Wilkins}


The polarization of periodically repeating systems is a discontinuous function of the atomic positions, a fact which seems at first to stymie attempts at their statistical learning. Two approaches to build models for bulk polarizations are compared: one in which a simple point charge model is used to preprocess the raw polarization to give a learning target that is a smooth function of atomic positions and the total polarization is learned as a sum of atom-centered dipoles and one in which instead the average position of Wannier centers around atoms is predicted. For a range of bulk aqueous systems, both of these methods perform perform comparatively well, with the former being slightly better but often requiring an extra effort to find a suitable point charge model. As a challenging test, we also analyze the performance of the models at the air–water interface. In this case, while the Wannier center approach delivers accurate predictions without further modifications, the preprocessing method requires augmentation with information from isolated water molecules to reach similar accuracy. Finally, we present a simple protocol to preprocess the polarizations in a data-driven way using a small number of derivatives calculated at a much lower level of theory, thus overcoming the need to find point charge models without appreciably increasing the computation cost. We believe that the training strategies presented here help the construction of accurate polarization models required for the study of the dielectric properties of realistic complex bulk systems and interfaces with ab initio accuracy.
